\section{Introducci\'on}

Este trabajo presenta un proyecto de Inteligencia de Negocios (BI) enfocado en 
analizar el transporte vehicular en la ciudad de Montevideo durante el año 2022
. El objetivo principal es procesar y unificar distintos conjuntos de datos, 
específicamente mediciones de velocidad promedio, conteo vehicular y 
siniestros de tránsito, para generar información útil. A través de este 
análisis, se busca entender cómo fluye el tráfico en la ciudad, evaluar la 
eficiencia de las avenidas y medir el peligro real de las distintas zonas 
relacionando los accidentes con el volumen de autos.

Para lograrlo, se armó un proceso de Extracción, Transformación y Carga (ETL) 
en Python que limpia los datos crudos y los organiza en un modelo de esquema 
estrella. Sobre esta base, se calculan métricas clave como el Índice de 
Eficiencia y el Índice de Riesgo, y se utilizan algoritmos como K-Means para 
clasificar los puntos críticos. Finalmente, todos los resultados y las 
recomendaciones, como la asignación de prioridades para instalar nuevos radares,
se muestran en reportes interactivos en Power BI para facilitar la lectura y 
la toma de decisiones. 




